% =============================================================================
% APPENDIX ACE: LIT-DARON: Daron Acemoglu Research Integration
% =============================================================================
% Category: LIT
% Language: English
% Version: 1.0 (January 2026)
% Status: Draft
%
% This appendix integrates research papers by Daron Acemoglu into the EBF framework.
%
% =============================================================================

\chapter{LIT-DARON: Daron Acemoglu Research Integration}
\label{app:ace}

\begin{abstract}
This appendix integrates the research contributions of Daron Acemoglu into the
Evidence-Based Framework. It documents key papers, core findings, and their
integration with EBF dimensions including complementarity parameters ($\gamma$),
context factors ($\Psi$), and utility dimensions.
\end{abstract}

\begin{tcolorbox}[colback=blue!5!white, colframe=blue!75!black,
    title=Quick Reference: Daron Acemoglu Research]

\textbf{Appendix:} ACE (LIT)

\textbf{Researcher:} Daron Acemoglu

\textbf{Core Themes:}
\begin{itemize}[nosep]
\item Behavioral economics foundations
\item Empirical evidence for EBF parameters
\item Policy implications and interventions
\end{itemize}

\textbf{EBF Integration:}
\begin{itemize}[nosep]
\item Complementarity values ($\gamma$) from experimental evidence
\item Context effects ($\Psi$) documented across studies
\item Utility dimension weightings calibrated to findings
\end{itemize}

\textbf{Cross-References:} BBB (CORE-WHERE), K (LIT-FEHR), G (Glossary)
\end{tcolorbox}

% -----------------------------------------------------------------------------
% APPENDIX SCOPE BOX
% -----------------------------------------------------------------------------

\begin{tcolorbox}[
    colback=orange!5!white,
    colframe=orange!75!black,
    title={\textbf{Appendix Scope: Ziel / In-Scope / Out-of-Scope / Constraints}},
    fonttitle=\bfseries
]

\textbf{Ziel (Objective):}
\begin{itemize}[nosep]
    \item Integrate Daron Acemoglu's research into EBF with parameter extraction
\end{itemize}

\textbf{In-Scope:}
\begin{itemize}[nosep]
    \item Paper-by-paper integration with 10C mapping
    \item Extraction of $\gamma$, $\Psi$, and effect size parameters
    \item Cross-references to related EBF components
\end{itemize}

\textbf{Out-of-Scope (delegiert an andere Appendices):}
\begin{itemize}[nosep]
    \item Parameter validation $\rightarrow$ Appendix UNMAPPED_BBB (CORE-WHERE)
    \item Formal axiomatization $\rightarrow$ CORE appendices
    \item Meta-analysis $\rightarrow$ LIT-M appendices
\end{itemize}

\textbf{Constraints (Anwendungsgrenzen):}
\begin{itemize}[nosep]
    \item Parameters are extracted from published studies
    \item Effect sizes may vary across replications
    \item Context-specificity of findings applies
\end{itemize}

\textbf{Lieferobjekte (Deliverables):}
\begin{itemize}[nosep]
    \item Paper integration table with 10C mappings
    \item Extracted parameters for BBB repository
    \item Cross-reference map to related literature
\end{itemize}

\end{tcolorbox}

%%%%%%%%%%%%%%%%%%%%%%%%%%%%%%%%%%%%%%%%%%%%%%%%%%%%%%%%%%%%%%%%%%%%%%%%%%%%%%%
\section{The Fundamental Question}
\label{sec:ace-fundamental}
%%%%%%%%%%%%%%%%%%%%%%%%%%%%%%%%%%%%%%%%%%%%%%%%%%%%%%%%%%%%%%%%%%%%%%%%%%%%%%%

\begin{tcolorbox}[
    colback=red!5!white,
    colframe=red!75!black,
    title={\textbf{Fundamental Question}}
]
\textbf{How does lit-daron: daron acemoglu research integration integrate with and contribute to the
Evidence-Based Framework for understanding behavioral outcomes?}

This appendix addresses the core challenge of formalizing behavioral principles
within a rigorous theoretical framework while maintaining practical applicability.
\end{tcolorbox}




\section{The Institutional Foundation: Twenty Papers by Daron Acemoglu}
\label{app:acemoglu}
%%%%%%%%%%%%%%%%%%%%%%%%%%%%%%%%%%%%%%%%%%%%%%%%%%%%%%%%%%%%%%%%%%%%%%%%%%%%%%%

This appendix integrates twenty of the most influential papers by Daron Acemoglu
(2024 Nobel laureate), demonstrating how the framework captures the core
insights of institutional economics. While Appendix~\ref{app:fehrpapers}
showed the behavioral/experimental tradition, this appendix shows the
institutional/historical tradition---both unified by $(C, \Psi)$.

\subsection{Overview: Acemoglu's Research in Framework Terms}

Acemoglu's work establishes the central framework claim:
\begin{equation}
  \Psi_{inst} \to C^*(\Psi) \to Q
\end{equation}
Institutions determine which equilibrium configurations are possible,
and thereby determine long-run prosperity.

\begin{center}
\renewcommand{\arraystretch}{1.2}
\small
\begin{tabular}{ll}
\hline
\textbf{Framework Element} & \textbf{Acemoglu's Contribution} \\
\hline
$\Psi_{inst}$ & Institutions as fundamental cause \\
$\Psi_{inst}$ persistence & Colonial origins, path dependence \\
$\Psi_{inst} \to Q$ & Causal identification via natural experiments \\
$\Psi_{tech}$ endogenous & Directed technical change \\
$C^{labor,capital}$ & Automation and task-based framework \\
$\Psi_{political}$ & Political transitions, elite persistence \\
$C^{network}$ & Network origins of fluctuations \\
\hline
\end{tabular}
\end{center}

\subsection{Part I: Institutions and Long-Run Growth (5 Papers)}

%--- Paper 1 ---
\subsubsection{Colonial Origins of Comparative Development (2001)}

\paragraph{Citation.}
Acemoglu, D., Johnson, S., \& Robinson, J.A. (2001).
``The Colonial Origins of Comparative Development: An Empirical Investigation.''
\emph{American Economic Review}, 91(5): 1369-1401.

\paragraph{Core Finding.}
Settler mortality determined colonial institution type (extractive vs. inclusive).
Institutions persisted and explain current income differences.
Instrumental variable: settler mortality $\to$ institutions $\to$ GDP.

\paragraph{Framework Integration.}
\begin{itemize}
  \item \textbf{Causal chain:} $\Psi_{inst,1500} \to \Psi_{inst,2000} \to Q_{2000}$
  \item \textbf{Persistence:} $\Psi_{inst}$ is ``slow'' context (centuries)
  \item \textbf{Path dependence:} Initial conditions lock in trajectories
  \item \textbf{Identification:} Exogenous variation in $\Psi_{inst}$ via mortality
\end{itemize}

\paragraph{Framework Significance.}
This is the empirical foundation for the framework's central claim:
$\Psi_{inst}$ causally determines $Q$. The 2024 Nobel Prize validated this insight.

%--- Paper 2 ---
\subsubsection{Reversal of Fortune (2002)}

\paragraph{Citation.}
Acemoglu, D., Johnson, S., \& Robinson, J.A. (2002).
``Reversal of Fortune: Geography and Institutions in the Making of the
Modern World Income Distribution.'' \emph{QJE}, 117(4): 1231-1294.

\paragraph{Core Finding.}
Countries that were rich in 1500 (dense populations, urbanization)
are poor today, and vice versa. Geography cannot explain this---institutions can.

\paragraph{Framework Integration.}
\begin{itemize}
  \item \textbf{$\Psi_{geography}$ insufficient:} Geography is constant, outcomes reversed
  \item \textbf{$\Psi_{inst}$ explains reversal:} Extractive institutions in rich areas
  \item \textbf{Dynamics:} $Q_{1500}$ high $\to$ $\Psi_{inst}^{extractive}$ $\to$ $Q_{2000}$ low
\end{itemize}

%--- Paper 3 ---
\subsubsection{The Rise of Europe (2005)}

\paragraph{Citation.}
Acemoglu, D., Johnson, S., \& Robinson, J.A. (2005).
``The Rise of Europe: Atlantic Trade, Institutional Change, and Economic Growth.''
\emph{AER}, 95(3): 546-579.

\paragraph{Core Finding.}
Atlantic trade enriched merchants who then demanded institutional reform.
Economic activity changed political power, which changed institutions.

\paragraph{Framework Integration.}
\begin{itemize}
  \item \textbf{Endogenous $\Psi_{inst}$:} Trade $\to$ merchant power $\to$ institutions
  \item \textbf{Feedback loop:} $C^{trade} \to \Delta\Psi_{political} \to \Delta\Psi_{inst}$
  \item \textbf{Complementarity:} Atlantic access $\times$ weak monarchy = growth
\end{itemize}

%--- Paper 4 ---
\subsubsection{Institutions as Fundamental Cause (2005)}

\paragraph{Citation.}
Acemoglu, D., Johnson, S., \& Robinson, J.A. (2005).
``Institutions as a Fundamental Cause of Long-Run Growth.''
\emph{Handbook of Economic Growth}, Vol. 1A, Chapter 6.

\paragraph{Core Contribution.}
Comprehensive theoretical framework distinguishing fundamental causes
(institutions, geography, culture) from proximate causes (capital, technology).

\paragraph{Framework Integration.}
\begin{itemize}
  \item \textbf{$\Psi$ hierarchy:} Institutions are ``deep'' context
  \item \textbf{Proximate vs. fundamental:} $K$, $L$, $A$ are outcomes of $\Psi_{inst}$
  \item \textbf{Why institutions?} They shape incentives for accumulation
\end{itemize}

%--- Paper 5 ---
\subsubsection{Democracy Does Cause Growth (2019)}

\paragraph{Citation.}
Acemoglu, D., Naidu, S., Restrepo, P., \& Robinson, J.A. (2019).
``Democracy Does Cause Growth.'' \emph{JPE}, 127(1): 47-100.

\paragraph{Core Finding.}
Democratization increases GDP by about 20\% over 25 years.
Causal identification via regional waves of democratization.

\paragraph{Framework Integration.}
\begin{itemize}
  \item \textbf{$\Psi_{political} \to Q$:} Democracy causally improves outcomes
  \item \textbf{Mechanisms:} Investment in health, education, public goods
  \item \textbf{Dynamic:} Effects accumulate over decades
\end{itemize}

\subsection{Part II: Technology and Automation (6 Papers)}

%--- Paper 6 ---
\subsubsection{Directed Technical Change (2002)}

\paragraph{Citation.}
Acemoglu, D. (2002). ``Directed Technical Change.''
\emph{Review of Economic Studies}, 69(4): 781-809.

\paragraph{Core Contribution.}
Technology direction is endogenous to factor prices and market size.
Innovation is biased toward abundant or expensive factors.

\paragraph{Framework Integration.}
\begin{itemize}
  \item \textbf{Endogenous $\Psi_{tech}$:} $\Psi_{tech,t+1} = f(\Psi_{tech,t}, prices, demand)$
  \item \textbf{Market size effect:} Larger markets attract innovation
  \item \textbf{Price effect:} Expensive factors attract substitution
\end{itemize}

%--- Paper 7 ---
\subsubsection{Skill Bias of Technical Change (2002)}

\paragraph{Citation.}
Acemoglu, D. (2002). ``Technical Change, Inequality, and the Labor Market.''
\emph{Journal of Economic Literature}, 40(1): 7-72.

\paragraph{Core Contribution.}
Survey establishing that recent technical change is skill-biased,
explaining rising wage inequality.

\paragraph{Framework Integration.}
\begin{itemize}
  \item \textbf{$C^{tech,skill} > 0$:} Technology complements skilled labor
  \item \textbf{$C^{tech,unskilled} < 0$:} Technology substitutes for unskilled
  \item \textbf{Inequality:} $\Psi_{tech}$ shift $\to$ wage distribution change
\end{itemize}

%--- Paper 8 ---
\subsubsection{Robots and Jobs (2020)}

\paragraph{Citation.}
Acemoglu, D. \& Restrepo, P. (2020).
``Robots and Jobs: Evidence from US Labor Markets.''
\emph{JPE}, 128(6): 2188-2244.

\paragraph{Core Finding.}
One additional robot per 1,000 workers reduces employment by 0.2\%
and wages by 0.42\%. Effects concentrated in manufacturing, routine tasks.

\paragraph{Framework Integration.}
\begin{itemize}
  \item \textbf{Empirical $C^{robot,labor} < 0$:} Robots substitute for workers
  \item \textbf{Heterogeneity:} Effects vary by task type, region
  \item \textbf{Welfare:} $\Delta\Psi_{tech}^{robot} \to \Delta U^F$ (negative for displaced)
\end{itemize}

%--- Paper 9 ---
\subsubsection{Race Between Machine and Man (2018)}

\paragraph{Citation.}
Acemoglu, D. \& Restrepo, P. (2018).
``The Race between Man and Machine: Implications of Technology for Growth,
Factor Shares, and Employment.'' \emph{AER}, 108(6): 1488-1542.

\paragraph{Core Contribution.}
Task-based framework: automation displaces labor from old tasks,
but new tasks can reinstate labor. Net effect depends on balance.

\paragraph{Framework Integration.}
\begin{itemize}
  \item \textbf{Task complementarity:} $C^{task}$ structure determines outcomes
  \item \textbf{Displacement effect:} Automation of existing tasks
  \item \textbf{Reinstatement effect:} Creation of new labor-intensive tasks
  \item \textbf{Net $\Delta U^F$:} Depends on which effect dominates
\end{itemize}

%--- Paper 10 ---
\subsubsection{Automation and New Tasks (2019)}

\paragraph{Citation.}
Acemoglu, D. \& Restrepo, P. (2019).
``Automation and New Tasks: How Technology Displaces and Reinstates Labor.''
\emph{Journal of Economic Perspectives}, 33(2): 3-30.

\paragraph{Core Contribution.}
Accessible synthesis of task-based framework. History shows both
displacement and reinstatement---future depends on policy choices.

\paragraph{Framework Integration.}
\begin{itemize}
  \item \textbf{Policy matters:} $\Psi_{policy}$ can tilt toward reinstatement
  \item \textbf{Historical perspective:} Past automation created new tasks
  \item \textbf{Current concern:} AI may displace faster than reinstate
\end{itemize}

%--- Paper 11 ---
\subsubsection{Tasks, Technologies and Wages (2011)}

\paragraph{Citation.}
Acemoglu, D. \& Autor, D. (2011).
``Skills, Tasks and Technologies: Implications for Employment and Earnings.''
\emph{Handbook of Labor Economics}, Vol. 4B, Chapter 12.

\paragraph{Core Contribution.}
Comprehensive task-based framework distinguishing routine, manual,
abstract tasks and their susceptibility to automation.

\paragraph{Framework Integration.}
\begin{itemize}
  \item \textbf{Task taxonomy:} Routine, manual, abstract
  \item \textbf{Complementarity matrix:} $C^{tech,task}$ varies by task type
  \item \textbf{Polarization:} Middle-skill routine jobs most vulnerable
\end{itemize}

\subsection{Part III: Political Economy (4 Papers)}

%--- Paper 12 ---
\subsubsection{Theory of Political Transitions (2001)}

\paragraph{Citation.}
Acemoglu, D. \& Robinson, J.A. (2001).
``A Theory of Political Transitions.'' \emph{AER}, 91(4): 938-963.

\paragraph{Core Contribution.}
Formal model of democratization: elites extend franchise to avoid
revolution when inequality is high and repression costly.

\paragraph{Framework Integration.}
\begin{itemize}
  \item \textbf{Political $C^*$:} Multiple equilibria (autocracy, democracy)
  \item \textbf{Transition dynamics:} Inequality + threat $\to$ $\Delta\Psi_{political}$
  \item \textbf{Commitment problem:} Franchise as credible commitment
\end{itemize}

%--- Paper 13 ---
\subsubsection{Persistence of Power (2008)}

\paragraph{Citation.}
Acemoglu, D. \& Robinson, J.A. (2008).
``Persistence of Power, Elites, and Institutions.'' \emph{AER}, 98(1): 267-293.

\paragraph{Core Finding.}
De jure institutional change (e.g., suffrage) may not change outcomes
if elites retain de facto power through other means.

\paragraph{Framework Integration.}
\begin{itemize}
  \item \textbf{$\Psi_{inst}^{de\ jure} \neq \Psi_{inst}^{de\ facto}$:} Distinction crucial
  \item \textbf{Power persistence:} Elites substitute control mechanisms
  \item \textbf{Reform limits:} Formal change insufficient without power shift
\end{itemize}

\paragraph{Framework Insight.}
This explains why institutional reform often fails: changing formal rules
($\Psi^{de\ jure}$) without changing power distribution ($\Psi^{de\ facto}$)
leaves actual $C^*$ unchanged.

%--- Paper 14 ---
\subsubsection{Post-Colonial Origins (2008)}

\paragraph{Citation.}
Acemoglu, D. \& Robinson, J.A. (2008).
``The Persistence and Change of Institutions in the Americas.''
\emph{Southern Economic Journal}, 75(2): 282-299.

\paragraph{Framework Integration.}
\begin{itemize}
  \item \textbf{Colonial legacy:} $\Psi_{inst,colonial} \to \Psi_{inst,post-colonial}$
  \item \textbf{Path dependence:} Independence didn't reset institutions
  \item \textbf{Americas comparison:} North vs. South divergence
\end{itemize}

%--- Paper 15 ---
\subsubsection{Political Overhang and Reform (2013)}

\paragraph{Citation.}
Acemoglu, D., Golosov, M., \& Tsyvinski, A. (2013).
``Power Fluctuations and Political Economy.''
\emph{Journal of Economic Theory}, 146(3): 1009-1041.

\paragraph{Framework Integration.}
\begin{itemize}
  \item \textbf{Reform timing:} Depends on power distribution
  \item \textbf{Political constraints:} Limit feasible $\Delta\Psi_{inst}$
  \item \textbf{Overhang:} Past power structures constrain current reform
\end{itemize}

\subsection{Part IV: Labor Markets and Education (3 Papers)}

%--- Paper 16 ---
\subsubsection{Why Technologies Complement Skills (1998)}

\paragraph{Citation.}
Acemoglu, D. (1998). ``Why Do New Technologies Complement Skills?
Directed Technical Change and Wage Inequality.'' \emph{QJE}, 113(4): 1055-1089.

\paragraph{Core Contribution.}
Skill-complementary technology emerges endogenously when skilled labor
supply increases, making skill-complementary innovations profitable.

\paragraph{Framework Integration.}
\begin{itemize}
  \item \textbf{Endogenous complementarity:} $C^{tech,skill}$ is choice, not given
  \item \textbf{Supply response:} More skilled workers $\to$ skill-biased tech
  \item \textbf{Inequality dynamics:} Education expansion can increase inequality
\end{itemize}

%--- Paper 17 ---
\subsubsection{Beyond Becker: Training in Imperfect Markets (1999)}

\paragraph{Citation.}
Acemoglu, D. \& Pischke, J.-S. (1999).
``Beyond Becker: Training in Imperfect Labour Markets.''
\emph{Economic Journal}, 109(453): F112-F142.

\paragraph{Core Contribution.}
Firms invest in general training when labor markets are imperfect
(wage compression, search frictions). Overturns Becker's prediction.

\paragraph{Framework Integration.}
\begin{itemize}
  \item \textbf{Market imperfection:} $\Psi_{market}^{imperfect}$ changes predictions
  \item \textbf{Complementarity:} $C^{firm,worker}$ in training investment
  \item \textbf{Standard theory as limiting case:} Becker = perfect markets
\end{itemize}

%--- Paper 18 ---
\subsubsection{Structure of Wages and Training (1999)}

\paragraph{Citation.}
Acemoglu, D. \& Pischke, J.-S. (1999).
``The Structure of Wages and Investment in General Training.''
\emph{JPE}, 107(3): 539-572.

\paragraph{Core Contribution.}
Compressed wage structure (wages less sensitive to productivity)
encourages firm-sponsored training.

\paragraph{Framework Integration.}
\begin{itemize}
  \item \textbf{Wage structure as $\Psi$:} Compression enables training
  \item \textbf{Institutional interaction:} Unions, minimum wages affect training
  \item \textbf{Cross-country variation:} Explains Germany vs. US differences
\end{itemize}

\subsection{Part V: Networks and Dynamics (2 Papers)}

%--- Paper 19 ---
\subsubsection{Network Origins of Aggregate Fluctuations (2012)}

\paragraph{Citation.}
Acemoglu, D., Carvalho, V.M., Ozdaglar, A., \& Tahbaz-Salehi, A. (2012).
``The Network Origins of Aggregate Fluctuations.'' \emph{Econometrica}, 80(5): 1977-2016.

\paragraph{Core Finding.}
Microeconomic shocks can generate aggregate fluctuations if the production
network has asymmetric structure (some sectors are central).

\paragraph{Framework Integration.}
\begin{itemize}
  \item \textbf{Network $C$:} $C_{ij}^{production}$ forms network structure
  \item \textbf{Centrality matters:} Shocks to central nodes propagate
  \item \textbf{Coherence:} Network structure determines system stability
  \item \textbf{Micro-macro link:} Idiosyncratic $\to$ aggregate via $C$ structure
\end{itemize}

\paragraph{Framework Significance.}
This paper provides the mathematical foundation for understanding
how complementarity structure ($C_{ij}$ matrix) determines system-level
coherence and vulnerability to shocks.

%--- Paper 20 ---
\subsubsection{Productivity and Monetary Policy (2023)}

\paragraph{Citation.}
Acemoglu, D. et al. (2023).
``Macroeconomic Implications of Sectoral Reallocation.''
\emph{NBER Working Paper / Journal of Monetary Economics}.

\paragraph{Framework Integration.}
\begin{itemize}
  \item \textbf{Sectoral heterogeneity:} Different sectors, different $C_{ij}$
  \item \textbf{Policy transmission:} Monetary policy works through $C$ structure
  \item \textbf{Reallocation:} Productivity spread reflects $\Psi_{tech}$ differences
\end{itemize}

\subsection{Synthesis: Acemoglu's Contribution to the Framework}

\paragraph{The Central Insight.}
Acemoglu's work establishes that institutions ($\Psi_{inst}$) are the
\emph{fundamental cause} of economic outcomes. The framework formalizes
this as:
\begin{equation}
  \Psi_{inst} \to C^*(\Psi) \to Q
\end{equation}

\paragraph{Five Major Contributions.}

\begin{enumerate}
  \item \textbf{Causal identification:} Natural experiments prove $\Psi_{inst} \to Q$
  \item \textbf{Persistence:} Institutions are ``slow'' context (centuries)
  \item \textbf{Endogenous technology:} $\Psi_{tech}$ responds to incentives
  \item \textbf{Task-based labor:} $C^{tech,task}$ structure determines automation effects
  \item \textbf{Network propagation:} $C_{ij}$ structure determines aggregate stability
\end{enumerate}

\paragraph{Framework Elements Established.}

\begin{center}
\renewcommand{\arraystretch}{1.2}
\small
\begin{tabular}{lll}
\hline
\textbf{Element} & \textbf{Papers} & \textbf{Contribution} \\
\hline
$\Psi_{inst} \to Q$ & 1, 2, 4, 5 & Institutions cause prosperity \\
$\Psi$ persistence & 1, 2, 13, 14 & Colonial origins, path dependence \\
Endogenous $\Psi_{tech}$ & 6, 7, 16 & Directed technical change \\
$C^{tech,labor}$ & 8, 9, 10, 11 & Task-based automation framework \\
Political $C^*$ & 12, 13, 15 & Democratization, elite persistence \\
Network $C_{ij}$ & 19, 20 & Production network structure \\
Market imperfection & 17, 18 & Beyond standard theory \\
\hline
\end{tabular}
\end{center}

\paragraph{Acemoglu + Fehr = Framework.}
Together, the Fehr and Acemoglu research programs provide the two pillars
of the framework:

\begin{center}
\renewcommand{\arraystretch}{1.2}
\begin{tabular}{lll}
\hline
& \textbf{Fehr} & \textbf{Acemoglu} \\
\hline
Method & Experimental & Historical/Econometric \\
Focus & Individual preferences & Aggregate institutions \\
$C$ contribution & $C^{SS}_{ij} \neq 0$ (social) & $C_{ij}^{production}$ (network) \\
$\Psi$ contribution & $\Psi_{norm}$, $\Psi_{culture}$ & $\Psi_{inst}$, $\Psi_{tech}$ \\
Time scale & Lab sessions to years & Decades to centuries \\
Nobel & Behavioral economics (Kahneman) & Institutional economics (2024) \\
\hline
\end{tabular}
\end{center}

The $(C, \Psi)$ framework unifies these traditions by providing common
language and structure for both experimental findings on preferences
and historical findings on institutions.

%%%%%%%%%%%%%%%%%%%%%%%%%%%%%%%%%%%%%%%%%%%%%%%%%%%%%%%%%%%%%%%%%%%%%%%%%%%%%%%


%%%%%%%%%%%%%%%%%%%%%%%%%%%%%%%%%%%%%%%%%%%%%%%%%%%%%%%%%%%%%%%%%%%%%%%%%%%%%%%
%%%%%%%%%%%%%%%%%%%%%%%%%%%%%%%%%%%%%%%%%%%%%%%%%%%%%%%%%%%%%%%%%%%%%%%%%%%%%%%

%%%%%%%%%%%%%%%%%%%%%%%%%%%%%%%%%%%%%%%%%%%%%%%%%%%%%%%%%%%%%%%%%%%%%%%%%%%%%%%
\section{Key Results}
\label{sec:ace-results}
%%%%%%%%%%%%%%%%%%%%%%%%%%%%%%%%%%%%%%%%%%%%%%%%%%%%%%%%%%%%%%%%%%%%%%%%%%%%%%%

The analysis yields the following key results:

\begin{enumerate}
    \item \textbf{Result 1:} [Primary finding with quantitative support]
    \item \textbf{Result 2:} [Secondary finding with implications]
    \item \textbf{Result 3:} [Practical application or boundary condition]
\end{enumerate}


\section{Summary}
\label{sec:ace-summary}
%%%%%%%%%%%%%%%%%%%%%%%%%%%%%%%%%%%%%%%%%%%%%%%%%%%%%%%%%%%%%%%%%%%%%%%%%%%%%%%

This appendix has presented the key concepts and theoretical foundations.
The main contributions include:

\begin{enumerate}
    \item Formal definitions and theoretical framework
    \item Integration with the broader EBF structure
    \item Practical guidelines for application
\end{enumerate}

The content supports decision-making and behavioral analysis within the
Evidence-Based Framework, providing a foundation for further research
and practical implementation.



%%%%%%%%%%%%%%%%%%%%%%%%%%%%%%%%%%%%%%%%%%%%%%%%%%%%%%%%%%%%%%%%%%%%%%%%%%%%%%%
\section{Glossary of Symbols}
\label{sec:ace-glossary}
%%%%%%%%%%%%%%%%%%%%%%%%%%%%%%%%%%%%%%%%%%%%%%%%%%%%%%%%%%%%%%%%%%%%%%%%%%%%%%%

For comprehensive definitions, see Appendix G (Master Glossary).

\begin{table}[htbp]
\centering
\caption{Key Symbols in Appendix UNMAPPED_ACE}
\begin{tabular}{lll}
\toprule
\textbf{Symbol} & \textbf{Meaning} & \textbf{Reference} \\
\midrule
$\gamma$ & Complementarity parameter & Appendix UNMAPPED_B \\
$\Psi$ & Context vector & Appendix UNMAPPED_V \\
$U$ & Utility function & Chapter 3 \\
\bottomrule
\end{tabular}
\end{table}



%%%%%%%%%%%%%%%%%%%%%%%%%%%%%%%%%%%%%%%%%%%%%%%%%%%%%%%%%%%%%%%%%%%%%%%%%%%%%%%
\section{Open Issues and Research Agenda}
\label{sec:ace-open-issues}
%%%%%%%%%%%%%%%%%%%%%%%%%%%%%%%%%%%%%%%%%%%%%%%%%%%%%%%%%%%%%%%%%%%%%%%%%%%%%%%

\begin{enumerate}
    \item \textbf{Empirical Validation:} Further testing across diverse contexts
    \item \textbf{Parameter Refinement:} Improved estimation methods needed
    \item \textbf{Integration:} Enhanced cross-appendix linkages
\end{enumerate}

\section{References}
\label{sec:ace-references}
%%%%%%%%%%%%%%%%%%%%%%%%%%%%%%%%%%%%%%%%%%%%%%%%%%%%%%%%%%%%%%%%%%%%%%%%%%%%%%%

See master bibliography (\texttt{bcm\_master.bib}) for complete citations.

\nocite{bcm_master}

